% !TeX root = ../../Book5.tex
% 纲要标题：### 14.4 结构应用
% 纲要位置：工作记录/计划纲要.md，第 4538 行。

\section{结构应用}
\label{b5:sec:1404}


% 纲要范围：
%
% * sl_2 三元组
% * 根弦的原型
% * 为一般根系和最高权理论提供模型
% * 区分环境矩阵代数与子代数中的三元组，按 sl_3 根弦核对零权重数与单链假设
%
% ---
%

三元组关系使 \(\mathfrak{sl}_2\) 可以嵌入较大的 Lie 代数，也使较大的表示在限制以后分成已经熟悉的权链。本节据此证明权链的整数性与反射对称性，为后面的根弦论证作准备。

\subsection{\texorpdfstring{\(\mathfrak{sl}_2\) 三元组}{sl2 三元组}}
\label{b5:subsec:1404:01}

\begin{definition}\label{b5:def:sl2-triple}
	设 \(\mathfrak g\) 为复 Lie 代数，\(x,y,h_0\in\mathfrak g\)。若它们满足
	\[
		[h_0,x]=2x,\qquad [h_0,y]=-2y,\qquad[x,y]=h_0
	\]
	则称它们为一个 \(\mathfrak{sl}_2\) 三元组。记矩阵单位为 \(E_{ij}\)，取 \(\mathfrak{sl}_2(\mathbb C)\) 的标准基 \(e=E_{12}\)、\(f=E_{21}\)、\(h=E_{11}-E_{22}\)。若三者不全为零，则由 \(e\mapsto x,f\mapsto y,h\mapsto h_0\) 给出的同态是单射，其像同构于 \(\mathfrak{sl}_2\)。
\end{definition}
上述赋值先按基 \(e,f,h\) 线性延拓；三条三元组关系恰好检查了全部基向量之间的括号，因而它确为 Lie 代数同态。其单射性来自\cref{b5:prop:sl2-simple}：同态的核为理想，而非零同态的核不能是整个单 Lie 代数。由单射性，非零三元组中的 \(x,y,h_0\) 均非零。
\begin{proposition}\label{b5:prop:nilpotent-matrix-sl2-triple}
	设 \(V\) 为有限维复向量空间，\(E\in\operatorname{End}_{\mathbb C}(V)\) 为幂零算子。则存在 \(F,H\in\operatorname{End}_{\mathbb C}(V)\)，使 \([H,E]=2E\)、\([H,F]=-2F\)、\([E,F]=H\)。
\end{proposition}
\begin{proof}
	对 \(E\) 取 Jordan 链，先处理长度 \(m+1\) 的一条，写成基 \(v_0,\ldots,v_m\)，使 \(Ev_i=v_{i-1}\)，其中 \(v_{-1}=0\)。定义
	\[
		Hv_i=(m-2i)v_i,\qquad Fv_i=(i+1)(m-i)v_{i+1},
	\]
	并令 \(v_{m+1}=0\)。代入可得 \([H,E]=2E\)、\([H,F]=-2F\)，以及
	\[
		[E,F]v_i=\bigl((i+1)(m-i)-i(m-i+1)\bigr)v_i=Hv_i.
	\]
	在全部 Jordan 链上取这些算子的直和，便得到所求 \(F,H\)。
\end{proof}
这里构造的是 \(\mathfrak{gl}(V)\) 中的三元组。如果 \(E\) 原本属于一个指定的真子代数，还须另外证明能否把 \(F,H\) 也选在该子代数中；本命题没有自动给出这个更强结论。

\begin{example}\label{b5:exm:upper-triangular-no-sl2-triple}
	在复 Lie 代数 \(\mathfrak{gl}_2(\mathbb C)\) 中，取 \(E=E_{12}\)。它是非零幂零矩阵，且
	\[
		F=E_{21},\qquad H=E_{11}-E_{22}
	\]
	与 \(E\) 构成三元组。但是在上三角子代数
	\[
		\mathfrak b_2=\left\{\,
		\begin{pmatrix}a&b\\0&c\end{pmatrix}
		\mid a,b,c\in\mathbb C\,\right\}
	\]
	中，不存在以 \(E\) 为第一个元素的三元组。事实上，对任意
	\(F'=\begin{pmatrix}a&b\\0&c\end{pmatrix}\in\mathfrak b_2\)，都有
	\[
		[E,F']=(c-a)E.
	\]
	若三元组存在，则第三个元素只能是 \(H'=[E,F']\)，从而 \([H',E]=0\)，与 \([H',E]=2E\ne0\) 矛盾。这里的障碍不只是标准选择 \(E_{21}\) 不在子代数中，而是子代数内任何 \(F'\) 都不能满足所需关系。
\end{example}

\subsection{根弦的原型}
\label{b5:subsec:1404:02}

\begin{proposition}\label{b5:prop:sl2-single-string}
	设 \(W\ne0\) 为有限维复 \(\mathfrak{sl}_2(\mathbb C)\) 模，\(d\in\mathbb C\) 为标准 Cartan 元 \(h=\operatorname{diag}(1,-1)\) 在 \(W\) 上的一个权。若全部权都在 \(d+2\mathbb Z\) 中，且每个权空间至多一维，则 \(W\) 不可约。存在非负整数 \(p,q\)，使权恰为
	\[
		d-2p,d-2p+2,\ldots,d+2q,\qquad d=p-q.
	\]
	对每个权 \(r\)，相反的权 \(-r\) 以相同重数出现。
\end{proposition}
\begin{proof}
	按\cref{b5:thm:sl2-complete-reducibility}，\(W\) 分解为若干 \(L_m\) 的直和。这些模型的权均为整数，所以 \(W\) 的全部权都是整数；特别地，作为其中一个权的 \(d\) 属于 \(\mathbb Z\)。由于所有权都在 \(d+2\mathbb Z\) 中，所有出现的最高权 \(m\) 都与 \(d\) 同奇偶。如果有两个直和项，那么偶数情形它们都含权零，奇数情形它们都含权一，均与权空间至多一维矛盾。因此只出现一个 \(L_m\)。其权无间断地从 \(m\) 降到 \(-m\)；以其中的 \(d\) 为原点记录两端，得到 \(d+2q=m\)、\(d-2p=-m\)，相加便是 \(d=p-q\)。对称性也由同一权链得到。
\end{proof}
这个判据的两个条件不能互相替代：一维性排除不了互不相遇的异奇偶链，同一奇偶性也排除不了多个直和项在某个权空间相遇。根空间中的应用还须区分关于整个对角子代数的权与关于单个 \(h\) 的权。

\begin{example}\label{b5:exm:sl3-adjoint-root-strings}
	在 \(\mathfrak g=\mathfrak{sl}_3(\mathbb C)\) 中，令 \(\mathfrak h\) 为迹零对角矩阵空间，并定义
	\[
		\varepsilon_i\bigl(\operatorname{diag}(t_1,t_2,t_3)\bigr)=t_i, \qquad \alpha=\varepsilon_1-\varepsilon_2,\qquad \beta=\varepsilon_2-\varepsilon_3.
	\]
	对 \(i\ne j\) 及 \(D\in\mathfrak h\)，有 \([D,E_{ij}]=(\varepsilon_i-\varepsilon_j)(D)E_{ij}\)。对 \(\gamma\in\mathfrak h^*\)，把满足 \([D,X]=\gamma(D)X\) 对全部 \(D\in\mathfrak h\) 成立的 \(X\in\mathfrak g\) 组成的空间记为 \(\mathfrak g_\gamma\)。由矩阵的各个坐标可得 \(\mathfrak g_{\varepsilon_i-\varepsilon_j}=\mathbb CE_{ij}\)（\(i\ne j\)），这些非零共同特征值泛函给出根方向；零泛函对应的空间则为 \(\mathfrak h\)。取三元组
	\[
		e_\alpha=E_{12},\qquad f_\alpha=E_{21},\qquad h_\alpha=E_{11}-E_{22}.
	\]
	由矩阵单位括号，得到
	\[
		\begin{aligned}
			[h_\alpha,E_{23}]&=-E_{23},& [h_\alpha,E_{13}]&=E_{13},\\
			[e_\alpha,E_{23}]&=E_{13},& [e_\alpha,E_{13}]&=0,\\
			[f_\alpha,E_{13}]&=E_{23},& [f_\alpha,E_{23}]&=0.
		\end{aligned}
	\]
	因此 \(W_\beta=\mathbb CE_{23}\oplus\mathbb CE_{13}\) 是三元组伴随作用下的子模；以 \(E_{13},E_{23}\) 为基，它的作用正是 \(L_1\) 的模型。这里升算子把根方向 \(\beta\) 送到 \(\beta+\alpha\)，降算子再送回；两端之外的作用均为零。以 \(\beta\) 为起点，初始权 \(d=\beta(h_\alpha)=-1\)，下降步数 \(p=0\)、上升步数 \(q=1\)，故
	\[
		\beta(h_\alpha)=p-q=-1.
	\]
	若以 \(\beta+\alpha\) 为起点，则 \(p=1,q=0\)，相应取值为 \(1\)。公式的符号由这两个实际升降方向决定。

	若改从 \(-\alpha\) 出发，沿 \(\alpha\) 方向经过零泛函。把各相应的共同特征空间全部取入，就得到
	\[
		W_0=\mathbb CE_{21}\oplus\mathfrak h\oplus\mathbb CE_{12}.
	\]
	为了找出 \(\mathfrak h\) 中与三元组交换的方向，写 \(H=\operatorname{diag}(a,b,c)\)、\(a+b+c=0\)。条件 \([H,E_{12}]=[H,E_{21}]=0\) 给出 \(a=b\)，再由迹为零得 \(c=-2a\)，所以这一方向由 \(K=\operatorname{diag}(1,1,-2)\) 张成。于是 \(\mathfrak h=\mathbb Ch_\alpha\oplus\mathbb CK\)，而 \(K\) 与三元组的三个元素都交换。三元组自身的线性生成空间在伴随作用下为 \(L_2\)，其最高权链为
	\[
		e_\alpha\ \xmapsto{\operatorname{ad}f_\alpha}\ -h_\alpha \ \xmapsto{\operatorname{ad}f_\alpha}\ -2f_\alpha \ \xmapsto{\operatorname{ad}f_\alpha}\ 0.
	\]
	于是有子模直和
	\[
		W_0=\operatorname{span}_{\mathbb C}\{e_\alpha,h_\alpha,f_\alpha\} \oplus\mathbb CK\cong L_2\oplus L_0.
	\]
	它关于 \(h_\alpha\) 的权为 \(2,0,-2\)，但零权空间为整个二维 \(\mathfrak h\)。两端的非零根空间虽然都是一维，单链判据要求的全部权空间一维却不成立。只由 \(E_{21}\) 生成的三元组子模是其中的 \(L_2\)；额外的平凡方向来自把整个 \(\mathfrak h\) 取入。
\end{example}

在更一般的根空间分解中应用单链判据，须先核对所取空间在三元组下稳定，再检查其中每个 \(h_\alpha\)-权空间至多一维。若 \(\alpha(h_\alpha)=2\)，则沿方向 \(\beta+k\alpha\) 的权为 \(\beta(h_\alpha)+2k\)，不同整数 \(k\) 给出互异的权。若 \(\alpha,\beta\) 线性无关，这些泛函均非零，所取空间的每个非零分量才都来自非零根空间。仅有 \(\beta\ne\pm\alpha\) 尚不足以排除其他成比例的情形；要作这样的简化，还需先证明关于成比例根的结论。若经过零泛函，则必须把相应的 Cartan 方向纳入，并重新核对权重数，不能只凭非零根空间一维应用判据。

\subsection{根系与最高权理论的模型}
\label{b5:subsec:1404:03}

设 \(\mathfrak g\) 为有限维复 Lie 代数，并在其中给定一个非零 \(\mathfrak{sl}_2\) 三元组。通过伴随作用，把整个 \(\mathfrak g\) 看成该 \(\mathfrak{sl}_2\) 的有限维模。于是 \(\operatorname{ad}h_0\) 可对角化，全部特征值为整数，\(\operatorname{ad}x\)、\(\operatorname{ad}y\) 则幂零。这里不要求整个 \(\mathfrak g\) 已经半单；使用的是限制到三元组后的表示定理。
\ProseParagraphBreak
这个结论为一般结构论提供了三件可用的事实：升降只把权移动固定的步长，有限链的端点受整数控制，而且链在反射 \(r\mapsto-r\) 下对称。在更大的半单 Lie 代数中，每一对相反根将产生一个这样的三元组；多个三元组的相容关系才进一步形成根系。至于不可约表示怎样由一组最高权整数共同决定，则需要后面的泛包络代数与三角分解。
